% SPDX-License-Identifier: CC-BY-SA-4.0
\documentclass[12pt]{article}

\usepackage{sbc-template}

% Workaround for sbc-template hook bug when hyperref is loaded
\AddToHook{begindocument/before}[sbc-template]{\relax}

\usepackage{graphicx,url}
\usepackage[utf8]{inputenc}
\usepackage[brazilian]{babel}
\usepackage{amsmath,amssymb}
\usepackage{booktabs}
\usepackage{csquotes}
\usepackage{newtxtext,newtxmath}
\usepackage{siunitx}
\sisetup{output-decimal-marker = {,}}
\usepackage{pgfplots}
\usepackage{pgfplotstable}
\pgfplotsset{compat=1.18}
\usepackage{microtype}
\usepackage{embedfile}
\usepackage[hidelinks]{hyperref}
\usepackage[brazilian,capitalize,nameinlink]{cleveref}

\input{data/values.tex}

\sloppy

\embedfile[
	desc={LaTeX source for this document},
	mimetype=text/plain
]{main.tex}

\embedfile[
	desc={BibTeX source file with references},
	mimetype=text/plain
]{references.bib}

\title{Engenharia de Algoritmos e Avaliação Microarquitetural\\de Variantes do Quicksort em C23}

\author{%
	Lucca M. A. Pellegrini,
	Ariel Inácio Jordão
}

\address{Instituto de Ciências Exatas e Informática\\
	Pontifícia Universidade Católica de Minas Gerais (PUC Minas)\\
	Belo Horizonte -- MG -- Brasil
	\email{\href{mailto:lucca@verticordia.com}{lucca@verticordia.com},
		\href{mailto:arielijordao@gmail.com}{arielijordao@gmail.com}}
}

\begin{document}

\maketitle

\begin{resumo}
	Este artigo apresenta uma análise experimental e microarquitetural exaustiva de variantes do algoritmo Quicksort implementadas no padrão C23. Avaliam-se estratégias de particionamento, seleção de pivô, híbridos com ordenações por inserção, salvaguardas adaptativas contra piores casos, execução iterativa e paralela em thread pool, e o impacto de compilação JIT e diretivas de laço sobre a microarquitetura AMD Zen 3.
\end{resumo}

\section{Introdução}\label{sec:intro}

O algoritmo Quicksort, proposto por Hoare~\cite{Hoare1961,Hoare1962}, é um dos pilares da ciência da computação. Alicerçado no paradigma de divisão e conquista, o método particiona recursivamente um vetor em torno de um elemento pivô, alcançando complexidade média assintótica em $\mathcal{O}(n \log n)$ com baixos fatores constantes~\cite{Sedgewick1978,Knuth1998}. Contudo, na microarquitetura contemporânea, o tempo de execução é amplamente governado pela interação com pipelines profundos, predição especulativa de desvios e hierarquias multinível de cache (L1, L2 e L3/LLC). Desvios imprevisíveis causam descartes de pipeline (\textit{flushes}) com penalidades de dezenas de ciclos, enquanto acessos sem localidade espacial provocam bloqueios de memória (\textit{stalls})~\cite{EdelkampWeiss2016}.

Ademais, o Quicksort clássico é suscetível a degeneração quadrática $\mathcal{O}(n^2)$ sob vetores pré-ordenados, invertidos, com chaves idênticas ou construídos por adversários determinísticos~\cite{McIlroy1999}. Para mitigar essas patologias, a literatura desenvolveu variantes com pivoteamento por mediana de três e \textit{ninther}~\cite{BentleyMcIlroy1993}, particionamento de três vias (3-way), hibridização com ordenação por inserção para subvetores pequenos~\cite{Sedgewick1978}, salvaguardas adaptativas como o PDQSort~\cite{Peters2021} e formulações iterativas ou paralelas.

Este trabalho realiza uma avaliação experimental e microarquitetural de variantes do Quicksort implementadas em C23 (\texttt{-std=c23}). Os objetivos compreendem: (i) calibrar empiricamente o limiar $M$ para métodos híbridos; (ii) mensurar a matriz canônica do projeto (Tradicional, Híbrido e Híbrido com Mediana de Três); (iii) confrontar o comportamento sob distribuições uniformes, quase ordenadas, repetidas e adversariais; (iv) avaliar extensões como Lomuto, Hoare, 3-Way, PDQSort, execução iterativa e paralelismo em \textit{thread pool}; e (v) correlacionar a contagem de operações com contadores de hardware (PMU) e a análise de código de máquina emitido via JIT.

\section{Fundamentação Teórica}\label{sec:teoria}

O Quicksort rearranja um arranjo $A[l..r]$ posicionando elementos em relação a um pivô $p$ e recursiona sobre as partições resultantes~\cite{Cormen2022}. O esquema de partição dita o fluxo de dados e os padrões de acesso:
\begin{itemize}
	\item \textbf{Lomuto:} Varredura unidirecional com ponteiro de fronteira~\cite{Cormen2022}. Embora didático, efetua $\mathcal{O}(n)$ trocas mesmo em vetores já ordenados e colapsa para $\mathcal{O}(n^2)$ sob chaves idênticas.
	\item \textbf{Hoare:} Dois ponteiros convergentes a partir dos extremos~\cite{Hoare1961,Hoare1962}. Minimiza trocas no caso médio e tolera duplicatas, mas não isola elementos iguais ao pivô.
	\item \textbf{3-Way (Bentley-McIlroy):} Segrega elementos em três faixas ($< p$, $= p$ e $> p$)~\cite{BentleyMcIlroy1993}. Evita recursões sobre chaves repetidas, atingindo tempo linear $\mathcal{O}(n)$ em arranjos com valores constantes.
\end{itemize}

A escolha do pivô condiciona o equilíbrio das partições. Pivôs fixos degradam sob entradas ordenadas. A \textbf{Mediana de Três} (Med3) adota a mediana entre o primeiro, o elemento central e o último elemento ($A[l], A[\lfloor (l+r)/2 \rfloor], A[r]$), neutralizando o pior caso em arranjos ordenados com custo fixo de três comparações~\cite{Sedgewick1978}. Para volumes extensos, o \textbf{Ninther de Tukey} calcula a mediana de três medianas de três elementos amostrados, aproximando a mediana real com custo reduzido~\cite{BentleyMcIlroy1993}.

A \textbf{hibridização} decorre do fato de que, para subvetores com $n \le M$, a sobrecarga recursiva e o particionamento superam o benefício assintótico do algoritmo~\cite{Sedgewick1978,Knuth1998}. A ordenação por inserção ($\mathcal{O}(M^2)$) torna-se substancialmente mais rápida nesses trechos devido a laços compactos e alta localidade espacial. Postergar essa ordenação e executar uma única varredura de inserção ao final do processo otimiza ainda mais o tempo global~\cite{Sedgewick1978}.

Em termos assintóticos, o Quicksort consome em média $2n \ln n \approx 1{,}386 n \log_2 n$ comparações~\cite{Sedgewick1978}. Partições assimétricas levam a $T(n) = T(n-1) + \mathcal{O}(n)$, resultando no pior caso $\mathcal{O}(n^2)$ e consumo de pilha linear, contornável via eliminação de recursão de cauda (\textit{tail-call}) para assegurar pilha em $\mathcal{O}(\log n)$~\cite{Sedgewick1978}. Na microarquitetura, o particionamento sobre chaves aleatórias exibe desvios condicionais com $\approx 50\%$ de probabilidade de tomada, degradando preditores de salto e gerando penalidades de falhas de predição (\textit{branch mispredictions})~\cite{EdelkampWeiss2016}, enquanto falhas de cache L1d e L2 reduzem o índice de instruções por ciclo (IPC).

\section{Trabalhos Relacionados}\label{sec:relacionados}

A evolução do Quicksort ilustra a transição da análise estritamente matemática para a engenharia de algoritmos sensível à arquitetura do processador. Hoare~\cite{Hoare1961,Hoare1962} introduziu a formulação original e o particionamento bidirecional. No presente estudo, o esquema de Hoare constitui a base do Quicksort Tradicional e serve de referência comparativa.

Sedgewick~\cite{Sedgewick1978} formalizou as diretrizes fundamentais da engenharia do Quicksort: o particionamento com mediana de três, o corte para ordenação por inserção em subvetores pequenos ($n \le M$) e a eliminação de recursão de cauda priorizando o menor subarranjo. Tais princípios balizam a variante Híbrida deste trabalho e estruturam o protocolo de calibração empírica de $M$.

Bentley e McIlroy~\cite{BentleyMcIlroy1993} desenvolveram a função \texttt{qsort} da biblioteca C da AT\&T, introduzindo particionamento ternário in-place por empacotamento de chaves iguais nas extremidades e seleção adaptativa entre mediana de três e ninther de Tukey. Essa variante integra nosso catálogo de extensões para avaliar o impacto de chaves repetidas e distribuições de baixa entropia.

Musser~\cite{Musser1997} introduziu o \textit{Introsort}, algoritmo que monitora a profundidade de recursão e comuta para o Heapsort ao ultrapassar $\lfloor 2 \log_2 n \rfloor$ níveis, limitando o pior caso estritamente a $\mathcal{O}(n \log n)$. A necessidade dessa proteção é demonstrada em nossos ensaios ao aplicar o adversário determinístico formulado por McIlroy~\cite{McIlroy1999}.

Peters~\cite{Peters2021} propôs o \textit{Pattern-Defeating Quicksort} (PDQSort), combinando detecção linear de arranjos ordenados ou decrescentes, particionamento com desvios previsíveis e salvaguarda contra o pior caso quadrático. O PDQSort representa em nosso arcabouço o padrão contemporâneo em algoritmos híbridos adaptativos.

Edelkamp e Weiß~\cite{EdelkampWeiss2016} demonstraram que falhas de predição de saltos constituem o principal gargalo do Quicksort moderno, concebendo o \textit{BlockQuicksort} para desvincular testes condicionais do fluxo de controle. Embora nossas implementações mantenham comparações diretas em C23, as conclusões de Edelkamp e Weiß fundamentam nossa telemetria de PMU (\textit{branch-misses} e IPC).

\section{Metodologia}\label{sec:metodologia}

O arcabouço experimental foi implementado em C23 (\texttt{-std=c23}), empregando recursos nativos como \texttt{<threads.h>}, \texttt{<stdatomic.h>}, \texttt{constexpr} e anotações \texttt{[[nodiscard]]}. A arquitetura isola o ruído observacional, garante reprodutibilidade determinística e extrai métricas de microarquitetura diretamente do processador.

\subsection{Arquitetura Dual-Pass}\label{sec:dual-pass}

A mensuração simultânea de operações algorítmicas (comparações e trocas) e grandezas microarquiteturais (tempo e ciclos de clock) introduz o efeito observador: incrementar contadores nos laços internos adiciona instruções e altera o escalonamento do compilador. Para eliminar essa distorção, o arcabouço adota um projeto em duas passagens:
\begin{enumerate}
	\item \textbf{Passada Limpa (\texttt{qs\_run\_clean}):} Executa a ordenação de forma pura sobre o vetor original, sem contadores internos. O tempo é medido em nanossegundos via \texttt{CLOCK\_MONOTONIC}, e a Unidade de Monitoramento de Desempenho (PMU) afere ciclos e eventos microarquiteturais.
	\item \textbf{Passada Contada (\texttt{qs\_run\_counted}):} Processa uma cópia idêntica do vetor original em versão instrumentada, registrando de forma exata e determinística o total de comparações, trocas (\textit{swaps}) e movimentações.
\end{enumerate}
Ao final de cada passada, a rotina \texttt{vec\_is\_sorted} audita a ordenação estrita e a integridade da soma de verificação multiconjunto das chaves.

\subsection{Calibração do Ponto de Corte \texorpdfstring{$M$}{M}}\label{sec:calibracao-m}

A eficiência da abordagem híbrida com inserção requer a calibração empírica do ponto de corte $M$. Valores excessivamente baixos sobrecarregam a recursão sobre subvetores diminutos; valores elevados expõem partes extensas do arranjo ao custo $\mathcal{O}(M^2)$ da inserção.

O protocolo experimental executa uma varredura paramétrica de $M \in [2, 128]$ sobre vetores pseudoaleatórios de dimensões $N = 50\,000$ e $N = 500\,000$, amostrando potências de dois e valores intermediários ($M \in \{2, 4, 8, 12, 16, 20, 24, 28, 32, 48, 64, 96, 128\}$). Para cada configuração, avalia-se o tempo da passada limpa, as instruções totais e o comportamento de cache.
Adotou-se $M^* = 50$, valor que minimiza o tempo médio em quatro das seis células de calibração (com 30 repetições por configuração). Entre $M = 30$ e $M = 64$ o tempo permanece em um platô cujas diferenças estão dentro do intervalo de confiança de 95\%, de modo que a escolha exata dentro desse platô tem impacto pequeno; $M^* = 50$ pertence ao conjunto sugerido pela especificação ($\{5, 10, 20, 30, 50\}$) e é adotado uniformemente em todas as variantes híbridas subsequentes.

\subsection{Cargas de Trabalho e Álgebra de Vetores}\label{sec:cargas}

Os experimentos cobrem um espectro abrangente de distribuições geradas a partir do PRNG determinístico xoshiro256**:
\begin{itemize}
	\item \textbf{Cargas Canônicas:} Vetores com distribuição uniforme (\texttt{random}), ordenados de forma crescente (\texttt{sorted}), decrescentes estritos (\texttt{reverse}), de valores repetidos com poucas chaves distintas (\texttt{repeated}) e degenerados para pivô fixo (\texttt{worst\_case}).
	\item \textbf{Cargas Estruturadas e Patológicas:} Vetores em tubo de órgão (\texttt{organ\_pipe}), dente de serra (\texttt{sawtooth}), vetores com fator contínuo de ordenação parcial $\alpha \in [0{,}0, 1{,}0]$ (\texttt{partial\_sort}), pior caso para mediana de três (\texttt{worst\_case\_med3}) e o vetor construído pelo adversário de McIlroy (\texttt{adversarial\_med3})~\cite{McIlroy1999}.
	\item \textbf{Álgebra de Vetores:} A linguagem de simulação DSL (\texttt{*.qs}) oferece operadores de composição: concatenação (\texttt{++}), intercalação (\texttt{**}), mistura estocástica (\texttt{\textasciicircum}), inversão (\texttt{\textasciitilde}) e fatiamento (\texttt{slice}).
\end{itemize}

O conjunto experimental contempla a matriz canônica do projeto (Tradicional, Híbrido e Híbrido com Mediana de Três) e variantes contemporâneas: Lomuto, Hoare, 3-Way (Bentley-McIlroy), PDQSort, Ninther de Tukey, execução iterativa em pilha explícita no heap, eliminação de cauda (\textit{tail-call}) e paralelismo por \textit{thread pool} com 4 e 6 threads de trabalho.

\subsection{Ambiente de Medição}\label{sec:ambiente}

As medições foram conduzidas no nó computacional caracterizado na \cref{tab:host-info}. O processador AMD Ryzen 5 5600X (microarquitetura Zen 3) dispõe de 6 núcleos físicos, 12 threads lógicas (SMT), caches privados L1d de 32 KiB (8 vias) e L2 de 512 KiB (8 vias) por núcleo, e uma cache L3 compartilhada de 32 MiB (32 vias).

\begin{table}[htbp]
\centering
\caption{Especificação do ambiente de medição e topologia do processador.}
\label{tab:host-info}
\IfFileExists{data/host_info.tex}{%
\input{data/host_info.tex}%
}{%
\begin{tabular}{ll}
\toprule
\textbf{Parâmetro} & \textbf{Especificação do Host} \\
\midrule
Processador & AMD Ryzen 5 5600X (6 núcleos físicos, 12 threads) \\
Hierarquia de Cache & L1d: 32 KiB, L2: 512 KiB (por núcleo); L3: 32 MiB unificada \\
Memória Principal & 31{,}28 GiB RAM \\
Sistema Operacional & Linux 6.18 LTS (\texttt{perf\_event\_paranoid=2}) \\
Compilador & Clang 23.1.4 (\texttt{-std=c23 -O3 -march=native}) \\
\bottomrule
\end{tabular}%
}
\end{table}

A telemetria acessa contadores de desempenho de hardware através de \texttt{perf\_event\_open}, registrando ciclos (\texttt{cycles}), instruções (\texttt{instructions}), saltos (\texttt{branches}), falhas de predição de saltos (\texttt{branch-misses}) e falhas de cache L1d (\texttt{L1-dcache-load-misses}). As contagens operam com privilégios isolados (\texttt{exclude\_kernel = 1} e \texttt{exclude\_hv = 1}), monitorando exclusivamente o processo em espaço de usuário.

\noindent\textbf{Ressalva Metodológica de Telemetria da PMU:} Na arquitetura AMD Zen 3 com \texttt{perf\_event\_paranoid=2}, o evento genérico \texttt{PERF\_COUNT\_HW\_CACHE\_LL} (\texttt{LLC-load-misses}) não é suportado pelo kernel sem privilégios elevados. Por essa razão, a grandeza de falhas de LLC (\texttt{mean\_llc\_misses}) é reportada formalmente como indisponível no ambiente de teste, evitando distorções decorrentes de contagens nulas.

\subsection{Protocolo de Execução e Isolamento}\label{sec:protocolo}

Para mitigar perturbações do sistema operacional:
\begin{enumerate}
	\item \textbf{Afinidade de CPU (\textit{Core Pinning}):} Cada thread de medição é fixada a um núcleo físico exclusivo via \texttt{sched\_setaffinity}, evitando migrações e preservando a afinidade de cache.
	\item \textbf{Controle de Estado de Cache (\textit{Cold Cache Enforcement}):} Em regime de cache frio, um buffer dedicado com duas vezes o tamanho da LLC ($2 \times \text{LLC} = 64\text{ MiB}$) é percorrido e sobrescrito antes de cada ordenação, expurgando linhas remanescentes de todos os níveis.
	\item \textbf{Repetições e Estatística:} Cada ponto experimental é submetido a $R$ repetições independentes ($R \ge 5$). Os resultados reportam média aritmética e desvio padrão amostral, descartando passagens de aquecimento.
\end{enumerate}

\subsection{Compilação Estática e Dinâmica JIT}\label{sec:compilacao}

O arcabouço opera sob modelo híbrido:
\begin{itemize}
	\item \textbf{Fast-Path Estático:} Na ausência de opções especiais, o kernel (\texttt{quicksort\_kernel.c}) é compilado com o binário principal em Clang com \texttt{-std=c23 -O3 -march=native}, permitindo otimizações interprocedurais e instrumentação para Otimização Guiada por Perfil (PGO via \texttt{make pgo}).
	\item \textbf{Módulo Dinâmico JIT:} Quando o arquivo de simulação especifica flags alternativas (\texttt{-O1}, \texttt{-O2}, \texttt{-O3}) ou diretivas \texttt{\#pragma} de laço (\texttt{\#pragma clang loop unroll} ou \texttt{\#pragma clang loop vectorize}), o sistema invoca o Clang sob demanda. O compilador produz uma biblioteca dinâmica em \texttt{build/jit/<hash>.so} e a respectiva listagem em assembly \texttt{build/jit/<hash>.s} (via \texttt{clang -S -fverbose-asm}), carregada em tempo de execução via \texttt{dlopen}/\texttt{dlsym} para correlação direta entre diretivas, código de máquina e telemetria da PMU.
\end{itemize}

\section{Resultados e Discussão}\label{sec:resultados}

Esta seção apresenta os resultados empíricos coletados em \ValTotalConfigs\ configurações de teste no ambiente caracterizado na \cref{tab:host-info}. Todas as execuções foram validadas quanto à integridade da ordenação (\texttt{sorted\_ok = \ValSortedOkVal} em \ValSortedOkPct\% dos ensaios) e submetidas a protocolo de isolamento com afinidade estática de núcleos (\textit{core pinning}) e expurgo total de cache ($\ValCachePurgeMultiplier\times\text{LLC} = \ValCachePurgeMib\text{ MiB}$) entre repetições.

\subsection{Calibração Empírica do Ponto de Corte}\label{sec:res-calibracao-m}

A \cref{fig:calibracao-m} ilustra o comportamento do tempo médio de execução em função do limiar de corte $M \in [\ValCalMinM, \ValCalMaxM]$ sobre as cargas \texttt{cal\_random}, \texttt{cal\_partial} (fator de ordenação $\ValCalPartialFactor$) e \texttt{cal\_repeated} (\ValCalRepeatedKeys\ chaves únicas) para dimensões $N = \ValCalNSmall$ e $N = \ValCalNLarge$ (\texttt{calibration\_m.csv}).

% SPDX-License-Identifier: CC-BY-SA-4.0
\pgfplotsset{
    discard if not/.style 2 args={
        filter discard warning=false,
        x filter/.append code={
            \edef\tempa{\thisrow{#1}}%
            \edef\tempb{#2}%
            \ifx\tempa\tempb
            \else
                \def\pgfmathresult{nan}%
            \fi
        }
    }
}
\begin{figure}[htbp]
\centering
\begin{minipage}{0.485\columnwidth}
\centering
\begin{tikzpicture}
\begin{axis}[
    width=\linewidth,
    height=5.4cm,
    xlabel={Limiar de corte $M$},
    ylabel={Tempo médio (\unit{\milli\second})},
    title={$N = 10\,000$},
    title style={font=\footnotesize\bfseries, yshift=-1.5ex},
    xmin=-2, xmax=66,
    xtick={0, 16, 30, 40, 50, 64},
    xticklabels={0, 16, 30, 40, 50, 64},
    tick label style={font=\tiny},
    label style={font=\scriptsize},
    grid=major,
    grid style={dashed, gray!30},
    legend to name=leg:calibracao-m,
    legend style={
        font=\footnotesize,
        legend columns=-1,
        draw=gray!40,
        fill=white,
        fill opacity=0.9,
        /tikz/every even column/.append style={column sep=1.2em}
    },
    unbounded coords=discard
]
\addplot[color=blue!80!black, mark=square*, mark size=1.0pt, semithick]
    table [
        col sep=comma,
        x=cutoff,
        y=mean_ms,
        discard if not={dataset}{cal_random},
        discard if not={n}{10000}
    ] {data/calibration_m.csv};
\addlegendentry{Aleatório}

\addplot[color=orange!85!black, mark=triangle*, mark size=1.1pt, semithick]
    table [
        col sep=comma,
        x=cutoff,
        y=mean_ms,
        discard if not={dataset}{cal_partial},
        discard if not={n}{10000}
    ] {data/calibration_m.csv};
\addlegendentry{Parcialmente ordenado}

\addplot[color=teal!80!black, mark=diamond*, mark size=1.1pt, semithick]
    table [
        col sep=comma,
        x=cutoff,
        y=mean_ms,
        discard if not={dataset}{cal_repeated},
        discard if not={n}{10000}
    ] {data/calibration_m.csv};
\addlegendentry{Chaves repetidas}

\draw[dashed, red!80!black, semithick] ({axis cs:50,0} |- {rel axis cs:0,0}) -- ({axis cs:50,0} |- {rel axis cs:0,1})
    node[pos=0.92, anchor=east, font=\scriptsize\bfseries, text=red!80!black] {$M^*=50$};
\end{axis}
\end{tikzpicture}
\end{minipage}\hfill
\begin{minipage}{0.485\columnwidth}
\centering
\begin{tikzpicture}
\begin{axis}[
    width=\linewidth,
    height=5.4cm,
    xlabel={Limiar de corte $M$},
    ylabel={Tempo médio (\unit{\milli\second})},
    title={$N = 50\,000$},
    title style={font=\footnotesize\bfseries, yshift=-1.5ex},
    xmin=-2, xmax=66,
    xtick={0, 16, 30, 40, 50, 64},
    xticklabels={0, 16, 30, 40, 50, 64},
    tick label style={font=\tiny},
    label style={font=\scriptsize},
    grid=major,
    grid style={dashed, gray!30},
    unbounded coords=discard
]
\addplot[color=blue!80!black, mark=square*, mark size=1.0pt, semithick]
    table [
        col sep=comma,
        x=cutoff,
        y=mean_ms,
        discard if not={dataset}{cal_random},
        discard if not={n}{50000}
    ] {data/calibration_m.csv};

\addplot[color=orange!85!black, mark=triangle*, mark size=1.1pt, semithick]
    table [
        col sep=comma,
        x=cutoff,
        y=mean_ms,
        discard if not={dataset}{cal_partial},
        discard if not={n}{50000}
    ] {data/calibration_m.csv};

\addplot[color=teal!80!black, mark=diamond*, mark size=1.1pt, semithick]
    table [
        col sep=comma,
        x=cutoff,
        y=mean_ms,
        discard if not={dataset}{cal_repeated},
        discard if not={n}{50000}
    ] {data/calibration_m.csv};

\draw[dashed, red!80!black, semithick] ({axis cs:50,0} |- {rel axis cs:0,0}) -- ({axis cs:50,0} |- {rel axis cs:0,1})
    node[pos=0.92, anchor=east, font=\scriptsize\bfseries, text=red!80!black] {$M^*=50$};
\end{axis}
\end{tikzpicture}
\end{minipage}\\[1.2ex]
\pgfplotslegendfromname{leg:calibracao-m}
\vspace{-0.5ex}
\caption{Calibração do limiar $M \in [0, 64]$ sobre arranjos aleatórios, com ordenação parcial e com repetições para $N \in \{10^4, 5\times 10^4\}$. O platô ótimo situa-se na faixa $M \in [30, 64]$, justificando a seleção uniforme de $M^* = 50$.}
\label{fig:calibracao-m}
\end{figure}


Em $M = \ValCalMZero$ (Quicksort puro, sem comutação para inserção na base), o algoritmo incorre na sobrecarga máxima de chamadas recursivas sobre subvetores unitários, consumindo $\ValCalRandMZeroMs\text{ ms}$ em $N = \ValCalNLarge$ sob dados aleatórios. A introdução de cortes pequenos ($M \in [\ValCalMCutSmallMin, \ValCalMCutSmallMax]$) atenua a profundidade recursiva, registrando $\ValCalRandMSixteenMs\text{ ms}$ em $M = \ValCalMSixteen$ (ganho de $\ValCalRandGainSixteenPct\%$).

À medida que $M$ alcança o intervalo $[\ValCalPlateauMin, \ValCalPlateauMax]$, consolida-se um platô de mínimo desempenho, registrando $\ValCalRandMFiftyMs\text{ ms}$ em $M = \ValCalMFifty$ (redução de $\ValCalRandGainFiftyPct\%$ frente a $M = \ValCalMZero$). A análise estatística sobre as \ValCalNumCells\ células fatoriais (\ValCalNumDists\ distribuições $\times$ \ValCalNumSizes\ tamanhos) atesta que $M = \ValCalMFifty$ é o mínimo estrito (\textit{argmin}) em \ValCalArgminFiftyCount\ das \ValCalNumCells\ configurações (\texttt{cal\_random} em \ValCalNSmallK\ e \ValCalNLargeK, \texttt{cal\_partial} em \ValCalNLargeK\ e \texttt{cal\_repeated} em \ValCalNSmallK). Nas duas células restantes (\texttt{cal\_partial} em \ValCalNSmallK\ e \texttt{cal\_repeated} em \ValCalNLargeK), onde $M=\ValCalMSixtyFour$ foi o mínimo numérico ($\ValCalPartMSixtyFourTenKMs\text{ ms}$ e $\ValCalRepMSixtyFourFiftyKMs\text{ ms}$), a diferença em relação a $M=\ValCalMFifty$ ($\ValCalPartMFiftyTenKMs\text{ ms}$ e $\ValCalRepMFiftyFiftyKMs\text{ ms}$) é de $\ValCalDiffPartTenKMs\text{ ms}$ e $\ValCalDiffRepFiftyKMs\text{ ms}$, valores que se situam rigorosamente dentro do intervalo de confiança de \ValCalConfPct\% ($\text{IC}_{\ValCalConfPct\%}$ com semi-larguras de $\pm \ValCalCiPartTenKMs\text{ ms}$ e $\pm \ValCalCiRepFiftyKMs\text{ ms}$). Sob ruído amostral, qualquer corte na faixa $[\ValCalPlateauMin, \ValCalPlateauMax]$ é estatisticamente indistinguível. Justifica-se, assim, a fixação de $M^* = \ValCalMFifty$ de maneira padronizada em todos os experimentos híbridos subsequentes.

\subsection{Matriz Canônica de Desempenho}\label{sec:res-matriz-tp}

A \cref{tab:tp-official-50k} consolida a matriz obrigatória do Trabalho Prático para $n = \ValCalNLarge$ (\texttt{tp\_official.csv}), confrontando as variantes Tradicional (\texttt{trad}, Lomuto com pivô fixo no último elemento e $M=\ValCalMZero$), Híbrido (\texttt{hybrid}, Lomuto com $M=\ValCalMFifty$) e Híbrido com Mediana de Três (\texttt{hybrid\_med3}, partição de Lomuto com mediana de três e $M=\ValCalMFifty$). A \cref{fig:scaling-tp} ilustra a escalabilidade temporal correspondente em escala log-log para $n \in \{\ValTpScalingNOne, \ValTpScalingNTwo, \ValTpScalingNThree, \ValTpScalingNFour\}$.

\begin{table}[htbp]
\centering
\scriptsize
\setlength{\tabcolsep}{2.2pt}
\caption{Matriz canônica de desempenho em $n = \ValCalNLarge$ (\texttt{tp\_official.csv}). A marcação $\dagger$ sinaliza estouro de profundidade recursiva ($d \ge \ValTpDepthLimit$) com chaveamento para Heapsort.}
\label{tab:tp-official-50k}
\begin{tabular}{@{}l rr rr rr l@{}}
\toprule
& \multicolumn{2}{c}{\textbf{Tradicional}} & \multicolumn{2}{c}{\textbf{Híbrido ($M=\ValCalMFifty$)}} & \multicolumn{2}{c}{\textbf{Híb.+Med3 ($M=\ValCalMFifty$)}} & \\
\cmidrule(lr){2-3}\cmidrule(lr){4-5}\cmidrule(lr){6-7}
\textbf{Carga} & \textbf{Tempo (ms)} & \textbf{Comparações} & \textbf{Tempo (ms)} & \textbf{Comparações} & \textbf{Tempo (ms)} & \textbf{Comparações} & \textbf{Estouro ($d\ge \ValTpDepthLimit$)} \\
\midrule
\texttt{random} & \ValTpRandTradMs & \ValTpRandTradCmp & \ValTpRandHybMs & \ValTpRandHybCmp & \ValTpRandMedMs & \ValTpRandMedCmp & Não \\
\texttt{sorted} & \ValTpSortTradMs & \ValTpSortTradCmp & \ValTpSortHybMs & \ValTpSortHybCmp & \ValTpSortMedMs & \ValTpSortMedCmp & Trad$^\dagger$, Hyb$^\dagger$ \\
\texttt{reverse} & \ValTpRevTradMs & \ValTpRevTradCmp & \ValTpRevHybMs & \ValTpRevHybCmp & \ValTpRevMedMs & \ValTpRevMedCmp & Trad$^\dagger$, Hyb$^\dagger$ \\
\texttt{repeated} & \ValTpRepTradMs & \ValTpRepTradCmp & \ValTpRepHybMs & \ValTpRepHybCmp & \ValTpRepMedMs & \ValTpRepMedCmp & Todos$^\dagger$ \\
\texttt{worst\_case} & \ValTpWorstTradMs & \ValTpWorstTradCmp & \ValTpWorstHybMs & \ValTpWorstHybCmp & \ValTpWorstMedMs & \ValTpWorstMedCmp & Trad$^\dagger$, Hyb$^\dagger$ \\
\texttt{worst\_case\_med3} & \ValTpWorstMedTradMs & \ValTpWorstMedTradCmp & \ValTpWorstMedHybMs & \ValTpWorstMedHybCmp & \ValTpWorstMedMedMs & \ValTpWorstMedMedCmp & Todos$^\dagger$ \\
\bottomrule
\end{tabular}
\end{table}

% SPDX-License-Identifier: CC-BY-SA-4.0
\pgfplotsset{
    discard if not/.style 2 args={
        filter discard warning=false,
        x filter/.append code={
            \edef\tempa{\thisrow{#1}}%
            \edef\tempb{#2}%
            \ifx\tempa\tempb
            \else
                \def\pgfmathresult{nan}%
            \fi
        }
    }
}
\begin{figure}[htbp]
\centering
\begin{tikzpicture}
\begin{axis}[
    width=0.88\linewidth,
    height=4.4cm,
    xmode=log,
    ymode=log,
    log basis x={10},
    log basis y={10},
    xlabel={Dimensão do vetor ($n$)},
    ylabel={Tempo médio (\unit{\milli\second})},
    xmin=800, xmax=65000,
    ymin=0.002, ymax=200,
    xtick={1000, 5000, 10000, 50000},
    xticklabels={$10^3$, $5\times 10^3$, $10^4$, $5\times 10^4$},
    tick label style={font=\tiny},
    label style={font=\scriptsize},
    grid=major,
    grid style={dashed, gray!30},
    legend pos=north west,
    legend style={font=\tiny, fill=white, fill opacity=0.85, draw=gray!40},
    unbounded coords=discard
]
\addplot[color=blue!80!black, mark=square*, mark size=1.1pt, semithick]
    table [
        col sep=comma,
        x=n,
        y=mean_ms,
        discard if not={dataset}{random},
        discard if not={variant}{trad}
    ] {data/tp_official.csv};
\addlegendentry{\texttt{random} -- Tradicional}

\addplot[color=cyan!80!black, mark=triangle*, mark size=1.2pt, semithick, dashed]
    table [
        col sep=comma,
        x=n,
        y=mean_ms,
        discard if not={dataset}{random},
        discard if not={variant}{hybrid}
    ] {data/tp_official.csv};
\addlegendentry{\texttt{random} -- Híbrido ($M=50$)}

\addplot[color=teal!80!black, mark=diamond*, mark size=1.2pt, semithick]
    table [
        col sep=comma,
        x=n,
        y=mean_ms,
        discard if not={dataset}{random},
        discard if not={variant}{hybrid_med3}
    ] {data/tp_official.csv};
\addlegendentry{\texttt{random} -- Híb.+Med3 ($M=50$)}

\addplot[color=red!80!black, mark=otimes*, mark size=1.2pt, semithick]
    table [
        col sep=comma,
        x=n,
        y=mean_ms,
        discard if not={dataset}{sorted},
        discard if not={variant}{trad}
    ] {data/tp_official.csv};
\addlegendentry{\texttt{sorted} -- Tradicional (Lomuto)}

\addplot[color=olive!80!black, mark=star, mark size=1.4pt, semithick]
    table [
        col sep=comma,
        x=n,
        y=mean_ms,
        discard if not={dataset}{sorted},
        discard if not={variant}{hybrid_med3}
    ] {data/tp_official.csv};
\addlegendentry{\texttt{sorted} -- Híb.+Med3 ($M=50$)}

\addplot[color=purple!80!black, mark=x, mark size=1.5pt, semithick, dashdotted]
    table [
        col sep=comma,
        x=n,
        y=mean_ms,
        discard if not={dataset}{worst_case_med3},
        discard if not={variant}{hybrid_med3}
    ] {data/tp_official.csv};
\addlegendentry{\texttt{worst\_med3} -- Híb.+Med3 ($M=50$)}

\node[font=\tiny\bfseries, text=red!80!black, anchor=west] at (axis cs:5500,10.5) {$\dagger$ Fallback Heapsort ($d \ge 4096$)};
\end{axis}
\end{tikzpicture}
\vspace{-1ex}
\caption{Escalabilidade assintótica (escala log-log) do tempo de execução em função de $n$. Sob vetores ordenados, o esquema tradicional (Lomuto com pivô fixo) degenera quadraticamente até a ativação da salvaguarda de introsort ($d \ge 4096$), enquanto a Mediana de Três preserva o regime ideal ($\approx 0{,}36\text{ ms}$ em $n=50\,000$).}
\label{fig:scaling-tp}
\end{figure}


Em vetores aleatórios (\texttt{random}), a hibridização proporciona redução consistente no tempo de execução: em $n=\ValCalNLarge$, a variante \texttt{hybrid} obtém $\ValTpRandHybMs\text{ ms}$ frente a $\ValTpRandTradMs\text{ ms}$ do \texttt{trad} (aceleração de $\ValTpRandHybSpeedupPct\%$). Essa economia decorre da eliminação da sobrecarga de chamadas de função e do processamento em laço denso de inserção, mesmo ocorrendo aumento de $\ValTpRandHybCmpIncPct\%$ nas comparações ($\ValTpRandHybCmp$ contra $\ValTpRandTradCmp$). A variante \texttt{hybrid\_med3} atinge $\ValTpRandMedMs\text{ ms}$ com $\ValTpRandMedCmp$ comparações, aproximando-se da cota teórica em $\ValTheoryTwo n \ln n \approx \ValTpTheoreticalCmpSci\times \ValTheoryTenToSix$.

\noindent\textbf{Ocorrência Mandatória de Estouro de Profundidade (\textit{depth\_exceeded}):} Um resultado fundamental reportado pelo arcabouço é o disparo da salvaguarda de recursão em exatamente \ValTpDepthExceededTotal\ configurações experimentais da matriz oficial (\texttt{tp\_official.csv}). As variantes da matriz adotam partição de Lomuto, isolando o impacto do pivô e do corte $M$. Sob distribuições ordenadas (\texttt{sorted}), invertidas (\texttt{reverse}) ou adversariais para pivô fixo (\texttt{worst\_case}), o pivô fixo no último elemento em \texttt{trad} e \texttt{hybrid} subtrai apenas um único elemento por partição, gerando profundidade recursiva linear $\Theta(n)$. Ao atingir o limite estático de segurança configurado na DSL (\texttt{max\_depth = \ValTpDepthLimit}), o mecanismo de introsort comuta deterministicamente para o Heapsort (\texttt{qs\_heapsort}), evitando o estouro da pilha de execução de \ValTpStackLimitMib\ MiB e preservando a corretude da ordenação.
O fenômeno ocorre em:
(i) \texttt{sorted}: $n \in \{\ValTpScalingNTwo, \ValTpScalingNThree, \ValTpScalingNFour\}$ para \texttt{trad} e \texttt{hybrid} (\ValTpSortedDepthCount\ linhas);
(ii) \texttt{reverse}: $n \in \{\ValTpScalingNTwo, \ValTpScalingNThree, \ValTpScalingNFour\}$ para \texttt{trad} e \texttt{hybrid} (\ValTpReverseDepthCount\ linhas);
(iii) \texttt{worst\_case}: $n \in \{\ValTpScalingNTwo, \ValTpScalingNThree, \ValTpScalingNFour\}$ para \texttt{trad} e \texttt{hybrid} (\ValTpWorstDepthCount\ linhas);
(iv) \texttt{repeated}: $n = \ValCalNLarge$ para \texttt{trad}, \texttt{hybrid} e \texttt{hybrid\_med3} (\ValTpRepeatedDepthCount\ linhas, pois o particionamento de Lomuto não isola chaves repetidas e degenera para profundidade linear mesmo com mediana de três); e
(v) \texttt{worst\_case\_med3}: $n = \ValCalNSmall$ para \texttt{trad} e \texttt{hybrid} (\ValTpWorstMedTenKDepthCount\ linhas) e $n = \ValCalNLarge$ para \texttt{trad}, \texttt{hybrid} e \texttt{hybrid\_med3} (\ValTpWorstMedFiftyKDepthCount\ linhas, totalizando \ValTpWorstMedDepthCount\ linhas).
Totalizam-se rigorosamente \ValTpDepthExceededTotal\ linhas com \texttt{depth\_exceeded = \ValSortedOkVal}. Em $n=\ValCalNLarge$ sob \texttt{sorted}, \texttt{trad} e \texttt{hybrid} demandam $\ValTpSortTradCmp$ comparações e $\approx \ValTpSortedTradMsApprox\text{ ms}$. Em contraste, \texttt{hybrid\_med3} seleciona a mediana exata no centro do vetor, completando a ordenação em apenas $\ValTpSortMedMs\text{ ms}$ com $\ValTpSortMedCmp$ comparações --- redução de $\ValTpSortMedCmpRedPct\%$ nas operações e ganho de $\ValTpSortMedSpeedupX\times$ em tempo.

\subsection{Sensibilidade a Piores Casos e Adversários}\label{sec:res-worst-case}

Embora a mediana de três neutralize o pior caso trivial de arranjos pré-ordenados, ela permanece suscetível a adversários determinísticos projetados contra a regra de três pontos. O experimento com o vetor assassino de McIlroy (\texttt{ext\_killer\_med3} em \texttt{extensions\_pdq.csv}) demonstra esse colapso: para $n = \ValCalNLarge$, a variante \texttt{hoare\_med3} sofre degradação patológica para $\approx n^2/\ValExtKillerTheoryDenom$, realizando $\ValExtKillerHoareComps$ comparações e despendendo $\ValExtKillerHoareMs\text{ ms}$ (ou $\ValExtKillerIterMs\text{ ms}$ em execução iterativa com pilha explícita).

Em contrapartida, o seletor \textit{ninther} de Tukey (\texttt{ninther\_hybrid}) e o algoritmo adaptativo PDQSort (\texttt{pdq\_sort}) neutralizam totalmente o ataque. Ao computar a mediana de três medianas de três elementos amostrados, o ninther quebra os empates simétricos orquestrados pelo adversário: o tempo de execução cai para $\ValExtKillerNintherMs\text{ ms}$ com $\ValExtKillerNintherComps$ comparações no \texttt{ninther\_hybrid}, e $\ValExtKillerPdqMs\text{ ms}$ com $\ValExtKillerPdqComps$ comparações no \texttt{pdq\_sort}. Ambas as variantes preservam o regime linear-logarítmico, operando mais de $\ValExtKillerSpeedupX\times$ mais rápido que a mediana de três.

Padrão análogo manifesta-se sob a distribuição em tubo de órgão (\texttt{ext\_organ\_pipe}): em $n=\ValCalNLarge$, \texttt{pdq\_sort} ($\ValExtOrganPdqMs\text{ ms}$, $\ValExtOrganPdqComps$ comparações) e \texttt{ninther\_hybrid} ($\ValExtOrganNintherMs\text{ ms}$, $\ValExtOrganNintherComps$ comparações) superam substancialmente o particionamento de Hoare clássico ($\ValExtOrganHoareMs\text{ ms}$, $\ValExtOrganHoareComps$ comparações) e o 3-Way não adaptativo ($\ValExtOrganThreeWayMs\text{ ms}$, $\ValExtOrganThreeWayComps$ comparações).

\subsection{Variantes Avançadas: 3-Way, PDQSort e Ninther}\label{sec:res-avancadas}

Sob cargas de chaves repetidas com poucas chaves distintas (\texttt{ext\_repeated} em $n = \ValCalNLarge$), o particionamento ternário de Bentley-McIlroy (\texttt{three\_way\_med3}) e o PDQSort (\texttt{pdq\_sort}) demonstram superioridade assintótica frente ao Quicksort clássico de duas vias. Enquanto \texttt{hoare\_med3} particiona elementos iguais em subintervalos subsequentes ($\ValExtRepHoareMs\text{ ms}$ e $\ValExtRepHoareComps$ comparações), o particionamento ternário agrega todas as chaves equivalentes ao pivô em um bloco central contíguo em uma única passada. Como consequência, \texttt{three\_way\_med3} e \texttt{pdq\_sort} reduzem as comparações para $\ValExtRepThreeWayComps$ e $\ValExtRepPdqComps$ respectivamente (redução de $\ValExtRepCmpRedRatio\times$), concluindo a ordenação em $\ValExtRepThreeWayMs\text{ ms}$ e $\ValExtRepPdqMs\text{ ms}$ (aceleração de $\ValExtRepSpeedupRatio\times$).

Em dados puramente aleatórios (\texttt{ext\_random}, $n = \ValCalNLarge$), a complexidade adicional das trocas e ponteiros do particionamento 3-Way acarreta sobrecarga moderada: \texttt{three\_way\_med3} consome $\ValExtRandThreeWayMs\text{ ms}$ e $\ValExtRandThreeWayComps$ comparações, frente a $\ValExtRandHoareMs\text{ ms}$ e $\ValExtRandHoareComps$ comparações do \texttt{hoare\_med3} (acréscimo de $\ValExtRandThreeWayOverheadPct\%$ no tempo). Já o PDQSort, por detectar adaptativamente a ausência de duplicatas, mantém-se competitivo em $\ValExtRandPdqMs\text{ ms}$ com apenas $\ValExtRandPdqComps$ comparações. Sob o vetor patológico de McIlroy (\texttt{ext\_killer\_med3}), como demonstrado, \texttt{pdq\_sort} ($\ValExtKillerPdqMs\text{ ms}$, $\ValExtKillerPdqComps$ comparações) e \texttt{ninther\_hybrid} ($\ValExtKillerNintherMs\text{ ms}$, $\ValExtKillerNintherComps$ comparações) superam o colapso quadrático de \texttt{hoare\_med3} ($\ValExtKillerHoareMsFull\text{ ms}$, $\ValExtKillerHoareComps$ comparações) em mais de $\ValExtKillerSpeedupX\times$. Em tubo de órgão (\texttt{ext\_organ\_pipe}), \texttt{pdq\_sort} ($\ValExtOrganPdqMs\text{ ms}$) e \texttt{ninther\_hybrid} ($\ValExtOrganNintherMs\text{ ms}$) superam \texttt{hoare\_med3} ($\ValExtOrganHoareMs\text{ ms}$) e \texttt{three\_way\_med3} ($\ValExtOrganThreeWayMs\text{ ms}$).

\subsection{Modelos de Execução: Pilha Explícita e Paralelismo}\label{sec:res-execucao}

O confronto entre a implementação recursiva com eliminação de chamada de cauda (\texttt{hoare\_med3}) e a execução iterativa com pilha explícita alocada em heap (\texttt{hoare\_iter}) comprova a maturidade do compilador Clang com `-O3`. Em todas as cargas e dimensões, os contadores de comparações ($\ValExtRandHoareIterComps$) e trocas ($\ValExtRandHoareSwaps$) são rigorosamente idênticos, e a variação no tempo de execução em $n=\ValCalNLarge$ sob \texttt{ext\_random} é de apenas $\ValExtRandIterDiffPct\%$ ($\ValExtRandHoareMs\text{ ms}$ contra $\ValExtRandIterMs\text{ ms}$), caracterizando equivalência operacional estrita e ausência de penalidade residual de chamada.

A avaliação do \textit{thread pool} C23 nativo com afinidade estática de núcleos (\texttt{sched\_setaffinity}) evidencia ganhos substanciais de vazão. Em $n = \ValCalNLarge$ sobre \texttt{ext\_random}, o paralelismo de tarefas reduz o tempo de $\ValExtRandHoareMs\text{ ms}$ (serial) para $\ValExtRandParallelFourMs\text{ ms}$ com \ValParThreadsFour\ threads de trabalho (\texttt{parallel\_4t}), alcançando \textit{speedup} de $\ValExtRandParallelFourSpeedup\times$ (eficiência paralela de $\ValExtRandParallelFourEff\%$). Com \ValParThreadsSix\ threads (\texttt{parallel\_6t}), o tempo foi de $\ValExtRandParallelSixMs\text{ ms}$ ($\ValExtRandParallelSixSpeedup\times$), refletindo a contenção no despachante de tarefas para matrizes de granularidade média. Contudo, ao expandir a dimensão para $n = \ValCalNHundredK$, a sobrecarga de sincronização dilui-se: \texttt{parallel\_4t} atinge $\ValExtRandHundredKParFourMs\text{ ms}$ ($\ValExtRandHundredKParFourSpeedup\times$) e \texttt{parallel\_6t} atinge $\ValExtRandHundredKParSixMs\text{ ms}$ ($\ValExtRandHundredKParSixSpeedup\times$ frente aos $\ValExtRandHundredKSerialMs\text{ ms}$ do kernel serial). Em vetores parcialmente ordenados (\texttt{ext\_partial\_80}, $n=\ValCalNLarge$), o tempo decresce de $\ValExtPartEightySerialMs\text{ ms}$ para $\ValExtPartEightyParFourMs\text{ ms}$ ($\ValExtPartEightyParFourSpeedup\times$ em \ValParThreadsFour\ threads), demonstrando a eficácia da partição balanceada concorrente.

\subsection{Telemetria Microarquitetural e Análise de Assembly}\label{sec:res-micro-jit}

A telemetria via PMU revela a estrita interdependência entre entropia dos dados, previsibilidade de saltos e instruções executadas por ciclo (IPC). Sob cargas aleatórias (\texttt{ext\_random}), as comparações de partição tomam caminhos com probabilidade próxima a $\ValExtRandPathProbPct\%$, acarretando taxas de erro de predição de desvios na faixa de $\ValExtRandBranchMissMinPct\%$ a $\ValExtRandBranchMissMaxPct\%$. Cada erro gera um descarte de pipeline de \ValZenDiscardCyclesMin\ a \ValZenDiscardCyclesMax\ ciclos no núcleo Zen~\ValZenArchGeneration, limitando o IPC a valores entre $\ValExtRandIpcMin$ e $\ValExtRandIpcMax$ (típico $\ValExtRandIpcTypMin$ a $\ValExtRandIpcTypMax$).

Em contrapartida, sob arranjos com regularidade (\texttt{ext\_organ\_pipe}, \texttt{ext\_repeated}, \texttt{ext\_killer\_med3} e \texttt{jit\_sorted}), os preditores de salto convergem rapidamente, reduzindo a taxa de falhas para menos de $\ValPmuRegularBranchMissCutoff\%$ ($\ValPmuKillerBranchMissPct\%$ no adversário de McIlroy, $\ValPmuSortedBranchMissPct\%$ no vetor ordenado e $\ValPmuOrganBranchMissPct\%$ no tubo de órgão). Sem a penalidade de descarte de pipeline, a microarquitetura do processador atinge taxas de IPC entre $\ValPmuRegularIpcMin$ e $\ValPmuRegularIpcMax$, operando próxima ao limite teórico de \ValPmuZenMaxDecoded\ instruções decodificadas por ciclo.

\noindent\textbf{Impacto da Otimização JIT e Diretivas de Laço (\texttt{jit\_assembly.csv}):} O estudo de compilação dinâmica sob demanda revela um contraexemplo notável ao senso comum relativo ao desenrolamento de laços. Em $n = \ValCalNHundredK$ sob dados aleatórios (\texttt{jit\_random}):
\begin{enumerate}
\item O baseline \texttt{cflags\_o3\_native} requer $\ValJitRandOThreeMs\text{ ms}$, consumindo $\ValJitRandOThreeInst$ instruções e $\ValJitRandOThreeCycles$ ciclos (IPC $\ValJitRandOThreeIpc$ e $\ValJitRandOThreeBranchMissPct\%$ de falhas de desvio);
\item A imposição de desenrolamento forçado via diretiva \texttt{\#pragma clang loop unroll\_count(\ValUnrollFactor)} (\texttt{pragma\_unroll\_4}) eleva o tempo para $\ValJitRandUnrollFourMs\text{ ms}$ ($+\ValJitRandUnrollFourMsIncPct\% \approx +\ValJitRandUnrollFourMsIncPctRound\%$) e amplia as instruções totais para $\ValJitRandUnrollFourInst$ ($+\ValJitRandUnrollFourInstIncPct\% \approx +\ValJitRandUnrollFourInstIncPctRound\%$).
\end{enumerate}
Embora o desenrolamento reduza ligeiramente a taxa de erro de salto (de $\ValJitRandOThreeBranchMissPct\%$ para $\ValJitRandUnrollFourBranchMissPct\%$), ele introduz código redundante de tratamento de resíduo (\textit{tail-loop}) e inflaciona a pegada de código, sobrecarregando a decodificação da CPU sem acelerar o avanço de ponteiros dependente de dados.
Ademais, a desativação de vetorização (\texttt{pragma\_no\_vec}, $\ValJitRandNoVecMs\text{ ms}$) gera alteração inferior a $\ValJitRandNoVecDiffThresholdPct\%$, confirmando que laços de partição com passos variáveis e saltos condicionais bidirecionais não admitem vetorização SIMD automática. A transição entre `-O1` ($\ValJitRandOOneMs\text{ ms}$), `-O2` ($\ValJitRandOTwoMs\text{ ms}$) e `-O3` ($\ValJitRandOThreeMs\text{ ms}$) apresenta variação inferior a $\ValJitRandOptVarPct\%$, indicando saturação de otimização escalar.
Ressalta-se que falhas de cache LLC (\texttt{mean\_llc\_misses}) permanecem nulas devido a restrições de permissão do kernel no contador genérico de hardware sob Zen~\ValZenArchGeneration, sendo omitidas da discussão conforme nota metodológica da \cref{sec:ambiente}.

\subsection{Discussão Integrada e Limitações}\label{sec:res-discussao}

O confronto com os modelos analíticos da literatura clássica confirma a aderência matemática do arcabouço:
\begin{itemize}
\item \textbf{Caso Médio:} A teoria de Hoare e Sedgewick prevê $C(n) \approx \ValTheoryTwo n \ln n \approx \ValTheoryLnFactor n \log_2 n$ comparações. Para $n = \ValCalNLarge$, a cota equivale a $\ValTheoryTwo \times \ValCalNLarge \ln(\ValCalNLarge) \approx \ValTheoryFiftyKCmp$ comparações. O resultado empírico de \texttt{hybrid\_med3} ($\ValTpRandMedCmp$ comparações) aproxima-se do valor analítico com discrepância de apenas $\ValTheoryDiscrepancyPct\%$.
\item \textbf{Pior Caso Quadrático:} Sob partição de Lomuto com pivô fixo, a teoria estabelece $n(n-\ValTheoryOne)/\ValTheoryTwo$ comparações. Para $n = \ValTheoryOneK$, $\ValTheoryOneK \times \ValTheoryNineNineNine / \ValTheoryTwo = \ValTheoryOneKCmp$ comparações, coincidindo rigorosamente com o total medido nas cargas \texttt{sorted} e \texttt{reverse}. Para $n \ge \ValTheoryFiveK$, o acionamento da guarda de introsort ($d \ge \ValTpDepthLimit$) trunca o colapso estritamente quadrático antes da saturação da pilha.
\item \textbf{Compensação da Hibridização:} A inserção em subvetores com $n \le \ValTheoryCutoffFifty$ substitui partições ineficientes por varreduras sequenciais de alto aproveitamento espacial, proporcionando ganho global de $\ValTheoryGainMinPct\text{--}\ValTheoryGainMaxPct\%$ no tempo a despeito do aumento moderado nas comparações.
\end{itemize}

\noindent\textbf{Limitações do Estudo:} Os ensaios foram conduzidos exclusivamente em um processador \ValHostCpuModel\ (microarquitetura Zen~\ValZenArchGeneration); os contadores de PMU operaram sob permissões de usuário não privilegiado (\texttt{perf\_event\_paranoid=\ValHostParanoidLevel}), impedindo a coleta direta de falhas de LLC; finalmente, os módulos JIT compilados dinamicamente foram carregados via \texttt{dlopen}, operando fora do alcance das otimizações interprocedurais e PGO do executável principal.

\section{Conclusão}\label{sec:conclusao}

Este trabalho apresentou uma investigação experimental e microarquitetural de variantes do algoritmo Quicksort sob o padrão C23. A adoção da arquitetura dual-pass permitiu quantificar de forma desacoplada o esforço algorítmico exato e os eventos de baixo nível de hardware, garantindo a reprodutibilidade dos ensaios em regime rigoroso de controle de estado de cache e afinidade estática de núcleos.

A calibração sistemática do limiar de corte evidenciou a formação de uma zona de ótimo estável para métodos híbridos com ordenação por inserção, justificando empiricamente a escolha de um ponto de corte robusto que maximiza o ganho de tempo sem sensibilidade a ruídos estocásticos de medição.

Os experimentos de estresse e adversários comprovaram as vulnerabilidades severas de esquemas canônicos de particionamento e pivoteamento fixo, ratificando a necessidade de salvaguardas adaptativas como o particionamento ternário de Bentley-McIlroy, a amostragem por ninther e mecanismos de introsort para assegurar robustez em ambientes de produção.

Por fim, a correlação entre telemetria de hardware e código gerado dinamicamente via JIT esclareceu os limites da otimização de laços no particionamento do Quicksort, demonstrando que a previsibilidade de desvios dita preponderantemente o índice de instruções por ciclo e que transformações como o desenrolamento de laços podem degradar o desempenho na presença de fluxo de controle dependente de dados.

\bibliographystyle{sbc}
\bibliography{references}

\appendix
\section{Apêndice: Tabelas Estendidas de Otimização JIT}\label{sec:apendice}

A \cref{tab:jit-assembly-full} apresenta os resultados completos obtidos na campanha de compilação dinâmica sob demanda (\texttt{jit\_assembly.csv}), avaliando níveis de otimização do Clang e diretivas \texttt{\#pragma} de laço sobre as cargas \texttt{jit\_random} e \texttt{jit\_sorted} nas dimensões $n \in \{\ValCalNSmall, \ValCalNLarge, \ValCalNHundredK\}$. O arquivo bruto encontra-se disponível no caminho relativo \texttt{report/data/jit\_assembly.csv}.

\begin{table}[htbp]
\centering
\scriptsize
\caption{Matriz experimental completa de compilação dinâmica JIT e telemetria PMU (\texttt{report/data/jit\_assembly.csv}).}
\label{tab:jit-assembly-full}
\begin{tabular}{lrrrrrr}
\toprule
\textbf{Variante JIT} & \textbf{$n$} & \textbf{Tempo (\unit{\milli\second})} & \textbf{Instruções} & \textbf{Ciclos} & \textbf{IPC} & \textbf{Falhas Desvio (\%)} \\
\midrule\multicolumn{7}{c}{\textbf{Carga: \texttt{jit\_random}}} \\\midrule
\texttt{-O1} & \ValCalTenKRaw & \ValJitTabRandTenKOOneMs & \ValJitTabRandTenKOOneInst & \ValJitTabRandTenKOOneCyc & \ValJitTabRandTenKOOneIpc & \ValJitTabRandTenKOOneBr\% \\
\texttt{-O2} & \ValCalTenKRaw & \ValJitTabRandTenKOTwoMs & \ValJitTabRandTenKOTwoInst & \ValJitTabRandTenKOTwoCyc & \ValJitTabRandTenKOTwoIpc & \ValJitTabRandTenKOTwoBr\% \\
\texttt{-O3 -march=native} (base) & \ValCalTenKRaw & \ValJitTabRandTenKOThreeMs & \ValJitTabRandTenKOThreeInst & \ValJitTabRandTenKOThreeCyc & \ValJitTabRandTenKOThreeIpc & \ValJitTabRandTenKOThreeBr\% \\
\texttt{pragma no\_unroll} & \ValCalTenKRaw & \ValJitTabRandTenKNoUnrollMs & \ValJitTabRandTenKNoUnrollInst & \ValJitTabRandTenKNoUnrollCyc & \ValJitTabRandTenKNoUnrollIpc & \ValJitTabRandTenKNoUnrollBr\% \\
\texttt{pragma unroll(\ValUnrollFactor)} & \ValCalTenKRaw & \ValJitTabRandTenKUnrollFourMs & \ValJitTabRandTenKUnrollFourInst & \ValJitTabRandTenKUnrollFourCyc & \ValJitTabRandTenKUnrollFourIpc & \ValJitTabRandTenKUnrollFourBr\% \\
\texttt{pragma no\_vec} & \ValCalTenKRaw & \ValJitTabRandTenKNoVecMs & \ValJitTabRandTenKNoVecInst & \ValJitTabRandTenKNoVecCyc & \ValJitTabRandTenKNoVecIpc & \ValJitTabRandTenKNoVecBr\% \\
\texttt{-O1} & \ValCalFiftyKRaw & \ValJitTabRandFiftyKOOneMs & \ValJitTabRandFiftyKOOneInst & \ValJitTabRandFiftyKOOneCyc & \ValJitTabRandFiftyKOOneIpc & \ValJitTabRandFiftyKOOneBr\% \\
\texttt{-O2} & \ValCalFiftyKRaw & \ValJitTabRandFiftyKOTwoMs & \ValJitTabRandFiftyKOTwoInst & \ValJitTabRandFiftyKOTwoCyc & \ValJitTabRandFiftyKOTwoIpc & \ValJitTabRandFiftyKOTwoBr\% \\
\texttt{-O3 -march=native} (base) & \ValCalFiftyKRaw & \ValJitTabRandFiftyKOThreeMs & \ValJitTabRandFiftyKOThreeInst & \ValJitTabRandFiftyKOThreeCyc & \ValJitTabRandFiftyKOThreeIpc & \ValJitTabRandFiftyKOThreeBr\% \\
\texttt{pragma no\_unroll} & \ValCalFiftyKRaw & \ValJitTabRandFiftyKNoUnrollMs & \ValJitTabRandFiftyKNoUnrollInst & \ValJitTabRandFiftyKNoUnrollCyc & \ValJitTabRandFiftyKNoUnrollIpc & \ValJitTabRandFiftyKNoUnrollBr\% \\
\texttt{pragma unroll(\ValUnrollFactor)} & \ValCalFiftyKRaw & \ValJitTabRandFiftyKUnrollFourMs & \ValJitTabRandFiftyKUnrollFourInst & \ValJitTabRandFiftyKUnrollFourCyc & \ValJitTabRandFiftyKUnrollFourIpc & \ValJitTabRandFiftyKUnrollFourBr\% \\
\texttt{pragma no\_vec} & \ValCalFiftyKRaw & \ValJitTabRandFiftyKNoVecMs & \ValJitTabRandFiftyKNoVecInst & \ValJitTabRandFiftyKNoVecCyc & \ValJitTabRandFiftyKNoVecIpc & \ValJitTabRandFiftyKNoVecBr\% \\
\texttt{-O1} & \ValCalHundredKRaw & \ValJitTabRandHundredKOOneMs & \ValJitTabRandHundredKOOneInst & \ValJitTabRandHundredKOOneCyc & \ValJitTabRandHundredKOOneIpc & \ValJitTabRandHundredKOOneBr\% \\
\texttt{-O2} & \ValCalHundredKRaw & \ValJitTabRandHundredKOTwoMs & \ValJitTabRandHundredKOTwoInst & \ValJitTabRandHundredKOTwoCyc & \ValJitTabRandHundredKOTwoIpc & \ValJitTabRandHundredKOTwoBr\% \\
\texttt{-O3 -march=native} (base) & \ValCalHundredKRaw & \ValJitTabRandHundredKOThreeMs & \ValJitTabRandHundredKOThreeInst & \ValJitTabRandHundredKOThreeCyc & \ValJitTabRandHundredKOThreeIpc & \ValJitTabRandHundredKOThreeBr\% \\
\texttt{pragma no\_unroll} & \ValCalHundredKRaw & \ValJitTabRandHundredKNoUnrollMs & \ValJitTabRandHundredKNoUnrollInst & \ValJitTabRandHundredKNoUnrollCyc & \ValJitTabRandHundredKNoUnrollIpc & \ValJitTabRandHundredKNoUnrollBr\% \\
\texttt{pragma unroll(\ValUnrollFactor)} & \ValCalHundredKRaw & \ValJitTabRandHundredKUnrollFourMs & \ValJitTabRandHundredKUnrollFourInst & \ValJitTabRandHundredKUnrollFourCyc & \ValJitTabRandHundredKUnrollFourIpc & \ValJitTabRandHundredKUnrollFourBr\% \\
\texttt{pragma no\_vec} & \ValCalHundredKRaw & \ValJitTabRandHundredKNoVecMs & \ValJitTabRandHundredKNoVecInst & \ValJitTabRandHundredKNoVecCyc & \ValJitTabRandHundredKNoVecIpc & \ValJitTabRandHundredKNoVecBr\% \\
\midrule\multicolumn{7}{c}{\textbf{Carga: \texttt{jit\_sorted}}} \\\midrule
\texttt{-O1} & \ValCalTenKRaw & \ValJitTabSortTenKOOneMs & \ValJitTabSortTenKOOneInst & \ValJitTabSortTenKOOneCyc & \ValJitTabSortTenKOOneIpc & \ValJitTabSortTenKOOneBr\% \\
\texttt{-O2} & \ValCalTenKRaw & \ValJitTabSortTenKOTwoMs & \ValJitTabSortTenKOTwoInst & \ValJitTabSortTenKOTwoCyc & \ValJitTabSortTenKOTwoIpc & \ValJitTabSortTenKOTwoBr\% \\
\texttt{-O3 -march=native} (base) & \ValCalTenKRaw & \ValJitTabSortTenKOThreeMs & \ValJitTabSortTenKOThreeInst & \ValJitTabSortTenKOThreeCyc & \ValJitTabSortTenKOThreeIpc & \ValJitTabSortTenKOThreeBr\% \\
\texttt{pragma no\_unroll} & \ValCalTenKRaw & \ValJitTabSortTenKNoUnrollMs & \ValJitTabSortTenKNoUnrollInst & \ValJitTabSortTenKNoUnrollCyc & \ValJitTabSortTenKNoUnrollIpc & \ValJitTabSortTenKNoUnrollBr\% \\
\texttt{pragma unroll(\ValUnrollFactor)} & \ValCalTenKRaw & \ValJitTabSortTenKUnrollFourMs & \ValJitTabSortTenKUnrollFourInst & \ValJitTabSortTenKUnrollFourCyc & \ValJitTabSortTenKUnrollFourIpc & \ValJitTabSortTenKUnrollFourBr\% \\
\texttt{pragma no\_vec} & \ValCalTenKRaw & \ValJitTabSortTenKNoVecMs & \ValJitTabSortTenKNoVecInst & \ValJitTabSortTenKNoVecCyc & \ValJitTabSortTenKNoVecIpc & \ValJitTabSortTenKNoVecBr\% \\
\texttt{-O1} & \ValCalFiftyKRaw & \ValJitTabSortFiftyKOOneMs & \ValJitTabSortFiftyKOOneInst & \ValJitTabSortFiftyKOOneCyc & \ValJitTabSortFiftyKOOneIpc & \ValJitTabSortFiftyKOOneBr\% \\
\texttt{-O2} & \ValCalFiftyKRaw & \ValJitTabSortFiftyKOTwoMs & \ValJitTabSortFiftyKOTwoInst & \ValJitTabSortFiftyKOTwoCyc & \ValJitTabSortFiftyKOTwoIpc & \ValJitTabSortFiftyKOTwoBr\% \\
\texttt{-O3 -march=native} (base) & \ValCalFiftyKRaw & \ValJitTabSortFiftyKOThreeMs & \ValJitTabSortFiftyKOThreeInst & \ValJitTabSortFiftyKOThreeCyc & \ValJitTabSortFiftyKOThreeIpc & \ValJitTabSortFiftyKOThreeBr\% \\
\texttt{pragma no\_unroll} & \ValCalFiftyKRaw & \ValJitTabSortFiftyKNoUnrollMs & \ValJitTabSortFiftyKNoUnrollInst & \ValJitTabSortFiftyKNoUnrollCyc & \ValJitTabSortFiftyKNoUnrollIpc & \ValJitTabSortFiftyKNoUnrollBr\% \\
\texttt{pragma unroll(\ValUnrollFactor)} & \ValCalFiftyKRaw & \ValJitTabSortFiftyKUnrollFourMs & \ValJitTabSortFiftyKUnrollFourInst & \ValJitTabSortFiftyKUnrollFourCyc & \ValJitTabSortFiftyKUnrollFourIpc & \ValJitTabSortFiftyKUnrollFourBr\% \\
\texttt{pragma no\_vec} & \ValCalFiftyKRaw & \ValJitTabSortFiftyKNoVecMs & \ValJitTabSortFiftyKNoVecInst & \ValJitTabSortFiftyKNoVecCyc & \ValJitTabSortFiftyKNoVecIpc & \ValJitTabSortFiftyKNoVecBr\% \\
\texttt{-O1} & \ValCalHundredKRaw & \ValJitTabSortHundredKOOneMs & \ValJitTabSortHundredKOOneInst & \ValJitTabSortHundredKOOneCyc & \ValJitTabSortHundredKOOneIpc & \ValJitTabSortHundredKOOneBr\% \\
\texttt{-O2} & \ValCalHundredKRaw & \ValJitTabSortHundredKOTwoMs & \ValJitTabSortHundredKOTwoInst & \ValJitTabSortHundredKOTwoCyc & \ValJitTabSortHundredKOTwoIpc & \ValJitTabSortHundredKOTwoBr\% \\
\texttt{-O3 -march=native} (base) & \ValCalHundredKRaw & \ValJitTabSortHundredKOThreeMs & \ValJitTabSortHundredKOThreeInst & \ValJitTabSortHundredKOThreeCyc & \ValJitTabSortHundredKOThreeIpc & \ValJitTabSortHundredKOThreeBr\% \\
\texttt{pragma no\_unroll} & \ValCalHundredKRaw & \ValJitTabSortHundredKNoUnrollMs & \ValJitTabSortHundredKNoUnrollInst & \ValJitTabSortHundredKNoUnrollCyc & \ValJitTabSortHundredKNoUnrollIpc & \ValJitTabSortHundredKNoUnrollBr\% \\
\texttt{pragma unroll(\ValUnrollFactor)} & \ValCalHundredKRaw & \ValJitTabSortHundredKUnrollFourMs & \ValJitTabSortHundredKUnrollFourInst & \ValJitTabSortHundredKUnrollFourCyc & \ValJitTabSortHundredKUnrollFourIpc & \ValJitTabSortHundredKUnrollFourBr\% \\
\texttt{pragma no\_vec} & \ValCalHundredKRaw & \ValJitTabSortHundredKNoVecMs & \ValJitTabSortHundredKNoVecInst & \ValJitTabSortHundredKNoVecCyc & \ValJitTabSortHundredKNoVecIpc & \ValJitTabSortHundredKNoVecBr\% \\
\bottomrule
\end{tabular}
\end{table}

\end{document}
